Kinetic models of the lipase-catalysed esterification of trimethylolpropane (TMP) with fatty acids, a route to fully bio-based biolubricants, must describe three competing Ping Pong bi bi reactions with few, correlated parameters. Nishimura et al. (\emph{Biomass} 2026, 6, 12) proposed a two-step semi-mechanistic method that first estimates the enzyme selectivities against the degree of reaction (``fingerprinting'') and then fixes them while fitting lumped constants against time. Because every data point of that study is public, we could put one chain of questions to identical data. (1)~Is the study reproducible? All \NCheckAll{} data checks and all fingerprint fits were reproduced exactly; several printed dynamic fits, however, were not at their optimum, the full model is structurally non-identifiable, and one case study could not be reproduced, so in-sample fit statistics alone are an unreliable guide to the model. (2)~What does the remaining misfit point to? Structured residuals indicate pseudo-processivity (re-acylation of an ester before it leaves the active site) in the mass-transfer-limited data and, in one profile, first-order enzyme deactivation. (3)~Can these mechanisms be added without losing the stabilising fingerprint information? PPK-FIP (pseudo-processive kinetics with a fingerprint-informed prior) treats the fingerprint estimate as a calibrated Bayesian prior rather than a fixed value, averages model forms by BIC, selects hyper-parameters by leave-one-case-out meta-validation, and uses an exact reparameterisation of the shared Ping Pong denominator that makes exhaustive validation cheap. (4)~Does it describe and predict better? PPK-FIP lowered the minimised sum of squares in all six fits (by \InAkSeventyRed--\InAkHundredRed\%) and, in leave-one-time-point-out cross-validation with both steps of the original method re-estimated in every fold, reduced the prediction error by \LotoAllRedOrig\% relative to the original method and by \LotoAllRedFull\% relative to its full variant (both $p<0.001$). Ablation traced the gain to pseudo-processivity and to the prior, with within-batch deactivation acting specifically in the 100\,°C profile. (5)~Are the mechanisms genuine? The processive composition path transferred across water activities, and Monte-Carlo experiments showed that the two-step method overestimates the selectivity $S_M$ by \SimSMBiasOrigPct\% when pseudo-processivity is present, whereas PPK-FIP is nearly unbiased and does not invent processivity when it is absent. (6)~Is mechanism still needed? Gaussian processes, MLP ensembles, the pretrained TabICL foundation model and neural ODEs interpolated densely sampled profiles but failed in extrapolation; a Gaussian-process discrepancy added to PPK-FIP (PPK-FIP-GP) gave the lowest interpolation error of all methods (\HybLotoRedOrig\% below the original method) while extrapolating like the mechanistic model. The fingerprint idea is therefore most valuable as prior information, pseudo-processivity should be modelled rather than absorbed into apparent selectivities, and kinetic parameters should be reported with out-of-sample validation.
